\subsection{CT14 --- Aggregate closure}\label{ct:14}\label{aggregate-closure-mass-bounds-on-residual-classes}
\ctsignature{\(P_{\ast\to14}\to \ctcert{C4};\ \ctint{\ctnone};\ \ctrec{P_{14\to9},P_{14\to10}}.\)}

CT14 compares a forced-mass lower bound with a class-capacity upper bound.
Its entry contract records the residual class, lower mass, capacity record,
carried-charge multiplicity, branch state, and C4 claim.  The core plan has
four nodes: scope, joint bounds certification, multiplicity classification,
and comparison certification.

\texttt{BoundsState} contains the lower and upper estimates but does not claim
that their units agree.  The multiplicity node either constructs an exact CT9
input for an unbounded member, an exact CT10 input for a missing finite label,
or a \texttt{MultiplicityState}.  Only the last state reaches the comparison
certificate and C4.  Thus a sublinear count cannot masquerade as a sublinear
carried capacity without the multiplicity proof.

\begin{itemize}
\tightlist
\item \textbf{S-Def} fixes the residual class and its branch provenance.
\item \textbf{S-Comp} is split between bounds, multiplicity, and final comparison.
\item \textbf{S-Rout} and \textbf{S-Trig} are exact CT9 and CT10 inputs.
\end{itemize}

\begin{figure}[p]
\centering
\resizebox{0.90\textwidth}{!}{%
\begin{tikzpicture}[x=1cm,y=1cm,every node/.style={font=\scriptsize}]
\node[bcwide] (entry) at (0,0) {{\tiny\ttfamily CT14.entry}\\\textbf{Validated residual class}};
\node[bcdec,bclemma] (scope) at (0,-2.2) {{\tiny\ttfamily CT14.decide.scope}\\\textbf{Finite capacity description?}};
\node[bcscopefail] (scopeout) at (-6.4,-2.2) {{\tiny\ttfamily CT14.terminal.scope}\\\textbf{Scope exit}};
\node[bcwide,bclemma] (bounds) at (0,-4.7) {{\tiny\ttfamily CT14.certify.bounds}\\\textbf{Lower mass and upper capacity}};
\node[bcroutechoice,bctheorem] (mult) at (0,-7.2) {{\tiny\ttfamily CT14.decide.\allowbreak multiplicity}\\\textbf{Carried multiplicity counted?}};
\node[bcroute,bcrouteneutral] (ct9) at (-6.4,-9.8) {{\tiny\ttfamily CT14.terminal.ct9}\\\textbf{CT 9 exit}};
\node[bcroute,bcrouteneutral] (ct10) at (6.4,-9.8) {{\tiny\ttfamily CT14.terminal.ct10}\\\textbf{CT 10 exit}};
\node[bcwide,bclemma] (compare) at (0,-9.8) {{\tiny\ttfamily CT14.certify.comparison}\\\textbf{Same-unit aggregate contradiction}};
\node[bcclosed] (c4) at (0,-12.3) {{\tiny\ttfamily CT14.terminal.c4}\\\textbf{C4}};
\draw[bcarr] (entry)--(scope); \draw[bcscopearr] (scope)--(scopeout);
\draw[bcarr] (scope)--(bounds); \draw[bcarr] (bounds)--(mult);
\draw[bchandoff] (mult)--node[bclabel,above,sloped]{unbounded}(ct9);
\draw[bchandoff] (mult)--node[bclabel,above,sloped]{missing label}(ct10);
\draw[bcarr] (mult)--node[bclabel,right]{counted}(compare); \draw[bcarr] (compare)--(c4);
\end{tikzpicture}%
}
\caption[Closure-tactic 14: exact aggregate machine]{CT14 keeps bound certification, multiplicity accounting, and the final comparison as separate proof boundaries.}
\label{fig:ct14-aggregate-closure}
\end{figure}
\FloatBarrier
